\subsubsection*{Quanto acrescenta a saída seguinte?}
\index{peso@peso!agregacao das saidas@agregação das saídas}
\index{prova de intervalo@prova de intervalo!dimensao e numero de saidas@dimensão e número de saídas}

As contagens anteriores permitem responder a uma pergunta de carteira.
Partimos de duas saídas. Se for preciso fazer um terceiro pagamento
na mesma transacção, quantos dados novos surgem do lado das provas?
E que diferença faria passar de três saídas para quatro?

Escrevamos \(p\) para o número efectivo de saídas e \(M\) para a
menor potência de dois que contém \(p\). A prova trabalha com
\(64M\) posições e vai reduzindo essa dimensão à metade. Cada
duplicação de \(M\) exige mais uma ronda, e cada ronda acrescenta
dois pontos de 32 bytes. No intervalo escolhido, os dois comprimentos
dos vectores da prova continuam a ocupar um byte cada um. O número
de provas agregadas também se conserva. Aplicando a fórmula aos
quatro corpos de prova, obtêm-se os comprimentos do quadro.
\[
\begin{array}{r|r|r|r|r}
p&M&64M&\text{rondas}&\text{corpo da prova (bytes)}\\ \hline
2 &2&128 &7 & 32(6+14)=640\\
3 &4&256 &8 & 32(6+16)=704\\
4 &4&256 &8 & 704\\
5 &8&512 &9 & 32(6+18)=768
\end{array}
\]
A terceira saída custa 64 bytes adicionais \emph{na prova}. A
quarta ocupa a capacidade da mesma dimensão e acrescenta zero bytes
nesta coluna. Ainda terá de deixar no prefixo a sua própria chave
e etiqueta, e na base RingCT o seu compromisso e o seu montante
cifrado. Também a verificação usa os dados de todas as saídas que
a declaração contém.

Podemos acabar a conta dos bytes com hipóteses bastante precisas.
Conservemos o número de entradas, o tamanho dos seus deslocamentos,
a mesma taxa e o mesmo comprimento de \texttt{extra}. Pensemos em
destinos com endereços principais, para os quais chega a mesma
chave de transacção. Uma nova saída traz um byte de montante visível
zero, um de tipo de alvo, 32 de chave e um de etiqueta. São
\(1+1+32+1=35\) no prefixo. Na base RingCT há ainda oito bytes
de montante cifrado e 32 de compromisso. Para as contagens da
tabela, o comprimento de \texttt{vout} cabe sempre num byte.
Assim, a despesa fixa por saída é \(35+8+32=75\) bytes. A
passagem de duas para três saídas aumenta a serialização em
\(75+64=139\) bytes, e a de três para quatro em 75.
Nesta comparação medimos comprimentos. Mudar os destinatários
obriga naturalmente a calcular de novo os valores das assinaturas
e dos compromissos.

Podemos agora apreciar a diferença entre este custo e o peso.
Para \(M=2\), a parcela de recuperação é zero. Para \(M=4\)
já vale
\[
 \left\lfloor\frac45(4\cdot320-704)\right\rfloor
       =\lfloor460{,}8\rfloor=460 .
\]
Logo, sob as mesmas hipóteses, acrescentar a terceira saída aumenta
o peso em \(139+460=599\), enquanto acrescentar a quarta o aumenta
em 75. O salto repete-se quando a quinta saída requer \(M=8\).
A parcela de recuperação passa então a
\(\lfloor4(2560-768)/5\rfloor=1433\), um acréscimo de 973.
A quinta saída acrescenta 139 bytes à serialização e
\(139+973=1112\) unidades ao peso. A partir do número de saídas,
o leitor consegue assim antecipar em que degraus muda a dimensão
da prova e em quais muda somente o número de objectos comprometidos.

Há dois números, pois, para responder à pergunta sobre a saída
seguinte. Os bytes contam o que precisa de ser transportado e
guardado na forma integral. O peso leva essa contagem à regra
segundo a qual a transacção ocupa espaço no bloco. Uma chave
adicional em \texttt{extra}, ou uma taxa cujo \emph{varint} cresça,
acrescenta os seus próprios bytes ao cálculo. É o campo
\texttt{extra} que precisamos agora de abrir.
